\documentclass[a4paper,10pt]{report}

\usepackage{xltxtra}
\usepackage{xgreek}
\usepackage{gfsneohellenic}
\usepackage{hyperref}
\usepackage[usenames,dvipsnames,svgnames,table]{xcolor}
\usepackage{graphicx}

\setromanfont[Mapping=tex-text]{GFS Neohellenic}
% \setsansfont[Mapping=tex-text]{DejaVu Sans}
% \setmonofont[Mapping=tex-text]{DejaVu Sans Mono}

\renewcommand{\thesection}{\Roman{section}} % Αρίθμηση ενοτήτων

\hypersetup{ % Ρύθμιση των χρωμάτων συνδέσμων
    colorlinks, linkcolor={cyan},
    citecolor={purple}, urlcolor={orange}
}

\title{
\vspace{-1cm}
\centering
\includegraphics[scale=0.8]{aginara-logo3.png}
\vspace{1cm}
\includegraphics[scale=0.6]{artichoke.jpg}
\vspace{0.5cm}
}
\author{Γιάννης Κατάκης \\Γιώργος Παναγόπουλος\\ Νεκταρία Ρέκκα\\Ανδρέας Γρίβας}
\date{\today}

\begin{document}
\maketitle
\tableofcontents

\newpage
\section{Ιδέα}
Η αγκινάρα είναι μια web εφαρμογή που έχει σκοπό να συμβουλεύει τον χρήστη οπτικά για την 
βέλτιστη επιλογή καταστήματος για μια λίστα
προιόντων που επιθυμεί να αποκτήσει. Η εφαρμογή μπορεί να συμβουλέψει τον χρήστη από άποψη: 
\begin{itemize}
 \item Τιμής
 \item Πληρότητας καλαθιού
\end{itemize}
Η Εφαρμογή μας χρησιμοποιεί το \href{http://www.data.gov.gr/dataset/17/}{dataset} με τις τιμές 
προιόντων καταναλωτών.\\

\section{Παράδειγμα Εκτέλεσης}

\begin{enumerate}
 \item Ο χρήστης εισάγει την διεύθυνση του
 \item Με βοήθεια από το autocomplete συμπληρώνει την λίστα προιόντων που τον ενδιαφέρουν
 \item Πατάει στο κουμπί search
 \item Ανάλογα με την επιλογή που έχει ενεργή επιστρέφονται τα καταστήματα που βρίσκονται σε ακτίνα
 4km και γίνεται οπτικοποίηση των βέλτιστων αποτελεσμάτων
 \item Με click πάνω στο κατάστημα εμφανίζονται περισσότερες πληροφορίες σχετικά με την επιλογή
\end{enumerate}

\section{Δεδομένα}
Χρησιμοποιώντας το \href{http://www.data.gov.gr/dataset/17/}{dataset} από το data.gov.gr 
φτιάξαμε μια βάση δεδομένων με τα παραπάνω δεδομένα. Ωστόσο πέρα από αυτό χρησιμοποιήσαμε
api googlemaps και προσθέσαμε έτσι στο dataset και την τοποθεσία με συντεταγμένες των 
καταστημάτων για όσα καταστήματα είχαν σωστή - ολοκληρωμένη διεύθυνση. 

\section{Επαναχρησιμοποίηση}
Στο \href{https://github.com/Pana-sonic/aginara/tree/master/Database}{git} υπάρχουν τα script 
για την δημιουργία της βάσης σύμφωνα με όσα dataset υπάρχουν στον φάκελο Datasets. Συνεπώς αρκεί
κανείς να πάρει τα csv αρχεία που έχουν τα δεδομένα και να τα βάλουν στον φάκελο Datasets και στην
συνέχεια να τρέξουν το reset.sh για να στηθεί η βάση και να έχουν πρόσβαση σε όλη την πρόσθετη
λειτουργικότητα. Επίσης θα μπορούσε κανείς μετά να φτιάξει ένα webservice για να δώσει τις 
συντεταγμένες των καταστημάτων προς τα έξω εφόσον αυτό είναι επιθυμητό.

\section{WebService API}

Τα webservice που στήσαμε δίνουν Rest api και επιστρέφουν json. Ακολουθεί περιγραφή των 
διαθέσιμων μεθόδων.
\subsection{Functions}

public String marketCarts(@FormParam("latitude") double latitude, @FormParam("longtitude") double longtitude, @FormParam("products") String products)
public String minList(@FormParam("latitude") double latitude, @FormParam("longtitude") double longtitude, @FormParam("products") String products)

\begin{verbatim}
HashMap <String marketName,Double distanceFromUser> superMAPket(String ip); 
ArrayList  <MarketResults> smartList(ArrayList <String> marketNames); 
TODO : 
HashMap <String marketName,Double distanceFromUser> superMAPket(String ip,String address); 

// to call threw ajax

HashMap <String productName,int id> findProduct(String substring);

Double getMinPrice(Product p); //looks for one product everywhere
Double getMinPrice(Product p,String location);//looks for one product in a certain location
Double getMinPrice(Product p,String marketNames); //looks for one product 
//everywhere at marketName markets
Double getMinPrice(Product p,String marketName, String location);
Double getMaxPrice(Product p); //looks for one product everywhere
Double getMaxPrice(Product p,String location); looks for one product in a certain location
Double getMaxPrice(Product p,String marketNames); //looks for one product everywhere 
//at marketName markets
Double getMaxPrice(Product p,String marketName, String location);
Double getAveragePrice(Product p); //looks for one product everywhere
Double getAveragePrice(Product p,String location); looks for one product in a certain location
Double getAveragePrice(Product p,String marketNames); //looks for one product everywhere 
//at marketName markets
Double getAveragePrice(Product p,String marketName, String location);

\end{verbatim}

\subsection{Classes}

\begin{verbatim}

public Market{
    private String marketName;
    private String address;
    private double latitude;
    private double longitude;
}

public Product{
    private int id;
    private String name;
    private double min,max,average;
}

public MarketResults{
    private Market market;
    private ArrayList <String> products;
    private double distance;
    private int percentage;
    private int actualSum; //sum of prices of products
    private int projectedSum; //expected sum of products
}

public PriceLocation{
    private String marketName;
    private double price;
}

\end{verbatim}

\section{Μελλοντική Δουλειά}

Στην συνέχεια η ιδέα είναι να εισαχθεί στην εφαρμογή δυνατότητα καταστημάτων 
που δεν υπάρχουν στο dataset να κάνουν εγγραφή στην εφαρμογή - δεν προλάβαμε να 
κάνουμε την υλοποίηση στα πλαίσια του hackathon. Θα δίνεται επιλογή στον χρήστη να 
επιβεβαιώσει πως είναι αποδεκτές οι τιμές για τη λίστα που συμπλήρωσε σε κάποιο 
κατάστημα και πως είναι πρόθυμος να αγοράσει. \\
Τότε θα ενημερώνονται όσα καταστήματα βρίσκονται στην ακτίνα αναζήτησης του χρήστη
και έχουν κάνει εγγραφή στην εφαρμογή και θα τους δίνεται η δυνατότητα να κάνουν
μια καλύτερη προσφορά στον χρήστη για την ίδια λίστα.\\
 

\end{document}
